% =====================================================================
%  Section: Introduction  (IJCAI 2027 revision: idea-centric)
% =====================================================================
\section{Introduction}
\label{sec:intro}

Classical paintings---whether oil on canvas, ink on silk, or pigment on
fresco---are cultural artifacts that encode centuries of human visual
imagination~\cite{lowenthal2015heritage}. From van Gogh's swirling
\textit{Starry Night} to Zhang Zeduan's panoramic \textit{Along the River
During the Qingming Festival}, these works imply motion, narrative, and
spatial depth that their static medium can only suggest. Animating such
paintings as explorable videos has therefore emerged as a compelling
research goal at the intersection of computer vision, generative
modeling, and digital cultural
heritage~\cite{zhang2024artgallery,dong2024enhancingpainting}.

Recent years have witnessed remarkable progress on two fronts that,
however, fail to meet on common ground. On one hand, image-to-video
(I2V) diffusion models~\cite{blattmann2023svd,wan2025,ltxvideo2024,
yang2025cogvideox}, combined with controllable conditioning
mechanisms~\cite{zhang2023controlnet,wang2024easycontrol,peng2024controlnext},
can animate a single photograph into a short clip with plausible
per-frame dynamics. Yet when applied directly to oil paintings or
ink-wash artworks, they suffer from two fundamental shortcomings.
First, they operate entirely in 2D pixel space: the model has no
geometric understanding of object depth ordering, self-occlusion, or
relative spatial layout, and consequently cannot synthesize novel
viewpoints or maintain consistent 3D structure across frames. Second,
long-range temporal coherence degrades rapidly beyond 2--4 seconds, as
the latent-space interpolation lacks explicit geometric constraints to
anchor the motion trajectory.

On the other hand, 3D reconstruction and rendering
methods~\cite{mildenhall2020nerf,kerbl20233dgs,huang20242dgs} can
reconstruct explorable 3D scenes from sparse views, and feed-forward
monocular variants~\cite{wang2024triposplat,tang2024dreamgaussian,
charatan2024pixelsplat} can even lift a single image to a 3D
representation. However, these reconstruction pipelines are optimized
for photorealistic fidelity; when applied to paintings, the rendering
process replaces expressive brushstrokes with geometric artifacts,
flattens the impasto texture, and bleeds semantic color boundaries
under novel-view projection.

The core tension is this: \emph{3D scene structure and artistic style
fidelity} are governed by incompatible optimization objectives, and no
single model family---neither pure diffusion nor pure 3D
reconstruction---satisfactorily achieves both. This tension becomes
especially acute when the animation task requires
\textbf{object-level manipulation}: moving a painted figure from one
compositional position to another while preserving its painterly
identity and filling the resulting gaps with plausible background
content.

\smallskip
\noindent\textbf{Our approach: modular separation of concerns.}
Rather than training a single monolithic network to jointly optimize
style and structure---an ill-posed objective whose gradients pull in
opposite directions---we propose to \emph{factor} the painting-to-video
problem into four interacting yet independently optimized sub-problems,
each addressed by a specialist and communicating with the next through
a small, interpretable handshake representation
(Figure~\ref{fig:overview}):

\begin{itemize}
    \item \textbf{PaintParse} handles \emph{perceptual grouping}:
    decomposing the input painting into foreground objects and a
    background layer using an art-domain-adapted segmentation backbone,
    augmented with monocular depth estimation for metric-scale layer
    ordering.
    \item \textbf{MonoSplat} handles \emph{geometric structure}:
    lifting each segmented foreground region into a high-fidelity 3D
    representation via feed-forward single-image 3D prediction,
    assembling an object-level model library with depth-derived metric
    rescaling.
    \item \textbf{SketchComp} handles \emph{compositional intent}:
    accepting hand-drawn line-art sketches specifying target object
    positions, solving a contour-based rigid registration to place each
    3D object onto its sketch, and filling disocclusion gaps through a
    depth-conditioned inpainting module that preserves background
    artistic texture.
    \item \textbf{ArtiFlow} handles \emph{temporal rendering}: rendering
    the composited 3D scene into multi-view keyframes and feeding them
    into a video diffusion backbone with dual control branches
    (depth + edge), producing temporally smooth animations with
    pixel-precise object motion governed by explicit 3D geometry
    rather than learned heuristics.
\end{itemize}

This decomposition enables a clean separation of concerns, and the
handshakes carry exactly the information each downstream stage needs
and no more, preventing error compounding. The result is a system that
supports \emph{arbitrary object repositioning}, \emph{multi-view
explorability}, and \emph{style-anchored animation}---capabilities
that no prior single-modality pipeline offers.

\smallskip
\noindent\textbf{Contributions.} Our contributions are:
\begin{enumerate}
    \item \textbf{A modular factorization} of the painting-to-video
    problem into four independently optimized specialists---PaintParse
    (perceptual grouping), MonoSplat (per-object 3D lifting),
    SketchComp (line-art composition with depth-aware inpainting), and
    ArtiFlow (dual-conditioned video synthesis)---communicating through
    interpretable handshake representations. To our knowledge this is
    the first end-to-end pipeline for 3D-aware, artistically
    controllable painting animation.
    \item \textbf{An object-level 3D lifting pipeline} (MonoSplat) that
    combines art-domain-adapted instance segmentation with feed-forward
    monocular 3D prediction and depth-derived metric rescaling,
    producing a reusable, retrievable library of 3D foreground objects
    from a single painting.
    \item \textbf{A sketch-conditioned 3D scene composition algorithm}
    (SketchComp) with a contour-based rigid registration that jointly
    optimizes region IoU and boundary Chamfer distance, and a
    depth-aware color inpainting module that fills disocclusion gaps
    while preserving the background's artistic texture.
    \item \textbf{A geometry-grounded video synthesis workflow}
    (ArtiFlow) with dual control branches (depth + edge) and
    first-frame latent anchoring for explicit 3D-structure-anchored
    animation with per-object motion control.
    \item \textbf{Comprehensive evaluation} on a curated benchmark of
    50 classical oil paintings and ink-wash works, achieving
    state-of-the-art style fidelity, perceptual quality, and temporal
    coherence, with user studies confirming strong preference for
    artistic authenticity, motion naturalness, and compositional-intent
    alignment.
\end{enumerate}
